\pagestyle{empty}

\section*{Repository használati útmutató}

\noindent
Az alábbi útmutató a dolgozathoz tartozó repository felépítését, valamint a reprodukálhatósághoz szükséges futtatási lépéseket foglalja össze. 
A projekt GitHubon érhető el: \texttt{https://github.com/RobertFerencsik/uni-thesis}

\subsection*{Beállítás}

\begin{enumerate}
    \item CUDA-képes videokártya szükséges a futtatáshoz.
    \item Python 3.10.0 telepítése.
    \item Csomagok telepítése: \texttt{pip install -r requirements.txt}
    \item A notebookok futtatásához a Jupyter környezet és az \texttt{ipykernel} csomag telepítése is szükséges.
\end{enumerate}

\subsection*{Használat}

\subsubsection*{Adathalmaz feldolgozása}

Az adathalmaz feldolgozásához a \texttt{notebooks/} mappában található notebookokat sorrendben szükséges futtatni.

\subsubsection*{Model}

A tároló gyökerében a \texttt{main.py} a belépési pont.

\paragraph{Hiperparaméter-keresés}

\begin{verbatim}
python main.py tune [--num-trials 20]
\end{verbatim}

\noindent
A próbák száma opcionális. A reprodukálhatóság érdekében az alapértelmezett érték \textbf{20}, ezzel lett mindhárom fázis futtatva.

\paragraph{Növekményes tanítóhalmazon tanulási görbe}

\begin{verbatim}
python main.py learning-curve [--num-portions 10]
\end{verbatim}

\noindent
A részek száma szintén opcionális. Az alapértelmezés \textbf{10}. Ezzel lett kiértékelve a növekményes adathalmazokon tanított modellek \textbf{általánosító} képessége.

\subsection*{Jegyzékfa}

Az alábbi fa segít eligazodni a tárolóba. A results \textbf{jegyzék} tartalmazza a futási eredményeket, így ennek megtekintésével futtatás nélkül is tanulmányozhatóak.

\begin{small}
\begin{verbatim}
uni-thesis/
|
+-- main.py                  # Belépési pont: tune vagy learning-curve
+-- requirements.txt         # pip-függőségek
+-- README.md
+-- .gitignore               # Nem verziókezelt fájlok listája
|
+-- notebooks/
|   +-- 1-eda-1.ipynb        # EDA
|   +-- 2-data-cleaning.ipynb
|   +-- 3-verification.ipynb
|   +-- 4-tokenization.ipynb # Tokenizáló
|   +-- 5-validation.ipynb   # Validációs lépések
|   `-- (*.png)              # Futtatáskor generált ábrák
|
+-- data/
|   +-- corpora/
|   |   +-- raw/             # Enron, CSV, train/teszt/validáció
|   |   `-- processed/       # Tisztított adatok, tokenizáló tanító szöveg
|   `-- models/              
|       +-- lstm_tokenizer.model
|       +-- train_ids.pkl, train_pieces.pkl
|       +-- learning_curve/
|       |   `-- portion_01..10/
|       |       +-- best_model.pt
|       |       +-- checkpoint_epoch_*.pt
|       |       +-- eval_test_metrics.json
|       |       +-- test_confusion_matrix.png
|       |       `-- training_curves.png
|       `-- tuning/          # learning_curve-vel egyező szerkezet
|
+-- src/
|   +-- config/              # config.json, hiperparaméterek
|   +-- data/                
|   +-- models/              # bilstm.py
|   +-- training/
|   +-- evaluation/
|   +-- experiments/         # tuning.py, learning_curve.py
|   `-- infrastructure/
|
+-- thesis/                  # Szakdolgozat LaTeX
|
`-- results/
    +-- preprocessing/       # Notebookok kimenete (Markdown, ábrák)
    +-- experiements/        # Modellek kiértékelései (archívum)
    |   +-- 1_phase_one_tuning/
    |   +-- 2_phase_two_learning_curve/
    |   +-- 2_phase_two_tuning/
    |   +-- 3_phase_three_learning_curve/
    |   `-- 3_phase_three_tuning/
    `-- SPAM_HAM_uzenetek_osztalyozasa_melytanulasi_modszerekkel.pdf
\end{verbatim}
\end{small}

